\documentclass[11pt]{article}

% 基础包
\usepackage[utf8]{inputenc}
\usepackage[T1]{fontenc}
\usepackage{amsmath, amssymb, amsfonts}
\usepackage{graphicx}
\usepackage{hyperref}
\usepackage{natbib}
\usepackage{geometry}
\geometry{margin=1in}

% 标题
\title{Memory-Efficient and Communication-Efficient Federated Learning}
\author{AutoSurvey Generated}
\date{\today}

\begin{document}

\maketitle

\section{Related Work}

\subsection{Foundations of Heterogeneous Federated Learning: Drift, Convergence, and Mitigation}
The foundational challenge in Federated Learning (FL) is the presence of non-IID client data, which induces \textit{client drift}—the divergence of local updates from the global objective. \citet{li2019convergence} provided the first rigorous convergence analysis of FedAvg under such heterogeneity and partial participation, establishing a critical trade-off between communication efficiency and convergence speed. This drift phenomenon is exacerbated by system heterogeneity, including unbalanced workloads \citep{daning2017weighted} and decentralized data distributions \citep{Tang2018D2DT}. To quantify heterogeneity, studies often employ Dirichlet-based data splits \citep{reguieg2023comparative}. A core theoretical advancement is the analysis by \citet{Yang2021AchievingLS}, which proved that linear speedup is achievable for FedAvg under non-IID data and partial worker participation. However, these analyses are often restricted to convex or strongly convex settings \citep{qu2023unified} and rely on bounded gradient assumptions. Momentum has been shown to provably benefit FL under unbounded heterogeneity \citep{Cheng2023MomentumBN}, and architectural strategies like widening networks in the Neural Tangent Kernel (NTK) regime can mitigate drift \citep{jian2025widening}. Yet, these solutions depend on idealized assumptions: momentum methods have limited scalability validation \citep{Cheng2023MomentumBN}, and NTK analyses require very wide networks \citep{jian2025widening}. Furthermore, non-uniform and time-varying client participation introduces bias that standard FedAvg cannot correct \citep{xiang2023towards}. While variance-reduction techniques like SCAFFOLD \citep{karimireddy2020scaffold} effectively correct drift, they incur communication overhead for maintaining control variates. Proximal methods for composite objectives \citep{Zhou2025FedCanonNC, zhang2025convex} offer convergence guarantees but often lead only to a neighborhood of the optimum. Non-parametric analyses in Reproducing Kernel Hilbert Spaces (RKHS) provide insights beyond stationary points \citep{su2023non}, but are limited to convex squared-loss regression. These limitations—assumptions of convexity, boundedness, full participation, and practical overhead—highlight a gap between theoretical drift mitigation and its application in resource-constrained systems.

\subsection{Communication-Efficient Federated Learning: Compression, Local Steps, and Decentralization}
To alleviate the communication bottleneck, three primary strategies have been developed: gradient compression, increased local training steps, and decentralized architectures. Compression techniques, including sparsification and quantization, reduce per-round data volume. Unified analyses \citep{Haddadpour2020FederatedLW} provide sharp convergence guarantees for FL combining local SGD with compression, but rely on assumptions of unbiased compressors and bounded variance. Error feedback and bias compensation mechanisms are essential for stabilizing training with biased compressors \citep{horvath2023stochastic}. However, their integration with FL under severe heterogeneity remains challenging, and they add memory overhead for storing error vectors. The paradigm of local training and periodic averaging, epitomized by FedAvg, directly trades computation for communication. Its convergence and representation learning capabilities have been analyzed \citep{li2019convergence, collins2022fedavg}, and accelerated variants like FedAc exist for convex settings \citep{Yuan2020FederatedAS}. Yet, these analyses often assume homogeneous data \citep{Yuan2020FederatedAS} or idealized participation models. Decentralized and peer-to-peer architectures eliminate the central server, using gossip or consensus protocols. Gradient Tracking (GT) methods are key for correcting bias in these settings \citep{koloskova2021improved}, and extensions support local steps \citep{Liu2023DecentralizedGT}. Momentum-augmented decentralized algorithms like Exact-Diffusion with Momentum (EDM) \citep{Hu2025ABD} and unified primal-dual frameworks for bilevel optimization \citep{Zhu2024SPARKLEAU} offer improved convergence. Asynchronous decentralized methods \citep{Lian2017AsynchronousDP} and gossip-based training \citep{blot2016gossip, Blot2018GoSGDDO} enhance scalability but rely on assumptions of bounded delays and doubly-stochastic mixing. Communication-efficient fully first-order methods for decentralized bilevel optimization \citep{Wen2024ACA} avoid second-order computations but have high complexity. Periodic global averaging can accelerate gossip SGD \citep{chen2021acceleraccelerating}. Lower bounds and accelerated algorithms exist for distributed optimization with compression \citep{He2023LowerBA}, and unbiased compression's impact on total communication cost has been formally analyzed \citep{He2023UnbiasedCS}. The primary limitations of these communication-efficient methods are their isolated analyses—they typically consider compression, local steps, or decentralization separately under restrictive assumptions (e.g., convexity, synchronous rounds, bounded heterogeneity). This isolation fails to capture the compounded effects when these techniques are combined, which is necessary for real-world deployments with joint constraints.

\subsection{Memory-Efficient Federated Learning: Subspace Methods, Low-Rank Optimization, and Implicit Regularization}
As model sizes grow, reducing the on-device memory footprint for optimizer states and gradients becomes critical. Core paradigms include subspace training and low-rank projection, which leverage the intrinsic low-rank structure of gradients. Gradient Low-Rank Projection (GaLore) \citep{Zhao2024GaLoreML} enables memory-efficient LLM training by projecting gradients onto a low-rank subspace. Federated adaptations like FedSub \citep{zhang2025efficient} constrain client updates to a low-dimensional subspace, reducing communication and memory while incorporating dual variables to mitigate drift. FedLoRU \citep{Park2024CommunicationEfficientFL} accumulates low-rank client updates to reconstruct a high-rank global model, providing communication efficiency and implicit regularization. However, these methods' theoretical guarantees rely on asymptotic random-matrix assumptions \citep{Park2024CommunicationEfficientFL} or are evaluated on limited scales \citep{zhang2025efficient}. Implicit regularization plays a key role, as gradient descent in constrained spaces can bias solutions toward better generalization \citep{Ma2017ImplicitRI, sekhari2021sgd}, and preconditioning can improve this effect \citep{su2024improving}. Advanced optimizer designs, such as applying Nesterov momentum to pseudo-gradients \citep{kallusky2025snoo} or matrix orthogonalization \citep{liu2025fedmuon}, aim to stabilize training but lack theoretical analysis under FL constraints. Integration with Parameter-Efficient Fine-Tuning (PEFT), specifically Low-Rank Adaptation (LoRA), is crucial for FL. \citet{Zhou2025ILoRAFL} addresses challenges like initialization-induced instability and rank incompatibility in federated LoRA, but adds algorithmic complexity. Zeroth-order (ZO) optimization, which uses only forward passes, offers radical memory savings. MeZO \citep{Malladi2023FineTuningLM} enables LLM fine-tuning with an inference-level memory footprint, and DeepZero \citep{chen2023deepzero} scales ZO methods for deep networks, but they require many more optimization steps. Randomized splitting algorithms for nonconvex composite FL \citep{Tran-Dinh2021FedDRR} handle nonsmooth regularizers but rely on standard smoothness assumptions. Extreme memory reduction techniques like ZeRO \citep{Rajbhandari2019ZeROMO} partition optimizer states but are designed for centralized training. The calibration of modern overparameterized networks, which can be poorly calibrated \citep{Guo2017OnCO}, is an additional concern when employing width-based strategies. The limitations of current memory-efficient FL methods are pronounced: their convergence analyses are often decoupled from communication constraints and data heterogeneity, their empirical validation is limited to moderate-scale benchmarks, and they do not jointly optimize for memory, communication, and robustness under extreme system heterogeneity. This creates a significant gap for deploying large models on resource-constrained, heterogeneous edge devices.

\bibliographystyle{plainnat}
\bibliography{references}

\end{document}
